% Chapter 1 File

\chapter{Introduction}
\label{introduction}
\thispagestyle{empty}


This is the introduction chapter of your thesis. Here you should:

\begin{itemize}
    \item Introduce the research topic and its context
    \item Explain the motivation for your research
    \item State the research problem or question
    \item Outline the objectives of your study
    \item Describe the structure of the thesis
\end{itemize}

\section{Background}

Provide background information that helps readers understand your research topic. This section should establish the context and importance of your work.

You can use \textbf{bold text}, \textit{italic text}, and \texttt{monospace text} for emphasis.

\section{Research Problem}

Clearly state the research problem or question that your thesis addresses. Explain why this problem is important and worth investigating.

\section{Research Objectives}

List the specific objectives of your research. For example:

\begin{enumerate}
    \item Objective 1: To investigate...
    \item Objective 2: To develop...
    \item Objective 3: To evaluate...
\end{enumerate}

\section{Thesis Structure}

Provide an overview of how the thesis is organized. For example:

Chapter 2 presents the literature review and theoretical framework. Chapter 3 describes the methodology used in this research. Chapter 4 presents the results and analysis. Finally, Chapter 5 concludes the thesis with a discussion of findings and future work.

\section{Example of Citations}

\ifusebibtex
You can cite references using the \texttt{cite} command. For example, according to
\cite{author2023}, this is an important finding. Multiple citations can be included
like this \cite{author2023, author2022}. The entries come from \texttt{References.bib}.
\else
This template is set to the manual bibliography (\texttt{useBibTeX} is \texttt{false}
in \texttt{config.tex}), so references are written out by hand in
\texttt{Bibliography.tex} and cited in the text by author and year, for example
``as shown by Author (2023)''.

Set \texttt{useBibTeX} to \texttt{true} to switch to \texttt{References.bib}. The
\texttt{\textbackslash cite} command then becomes available and this paragraph is
replaced by a worked example, without any other change to the document.
\fi

\section{Example of Abbreviations and Symbols}

Entries defined in \texttt{Lists.tex} only show up in the List of Abbreviations
and the List of Symbols once they are actually used in the text. The first use of
an acronym is expanded automatically, \gls{ai}, and every later use is short:
\gls{ai}. The same applies to \gls{ml} and \gls{api}.

Symbols behave the same way: \gls{symb:Pi} and \gls{symb:Lambda}.

\section{Example of Figures}

You can include figures in your thesis. Here's an example:

\begin{figure}[h]
    \centering
    % Uncomment and modify the following line when you have an actual figure
    % \includegraphics[width=0.8\linewidth]{figures/example_figure.png}
    \caption{Example figure caption. Always provide descriptive captions for your figures.}
    \label{fig:example}
\end{figure}

You can reference figures using \texttt{ref} command: see Figure \ref{fig:example}.

For side-by-side panels use the \texttt{subfigure} environment from
\texttt{subcaption}. Each panel needs an explicit width and carries its own
\texttt{caption}; \texttt{ref} then yields `1.1a', `1.1b' and so on.

% \begin{figure}[h]
%     \centering
%     \begin{subfigure}{0.45\linewidth}
%         \centering
%         \includegraphics[width=\linewidth]{figures/left.png}
%         \caption{Side view}
%         \label{fig:side}
%     \end{subfigure}\hfill
%     \begin{subfigure}{0.45\linewidth}
%         \centering
%         \includegraphics[width=\linewidth]{figures/right.png}
%         \caption{Front view}
%         \label{fig:front}
%     \end{subfigure}
%     \caption{The device seen from two angles}
%     \label{fig:device}
% \end{figure}

\section{Example of Tables}

Tables can be included as follows:

\begin{table}[h]
\centering
\caption{Example table showing sample data}
\label{tab:example}
\begin{tabular}{|l|c|r|}
\hline
\textbf{Column 1} & \textbf{Column 2} & \textbf{Column 3} \\
\hline
Row 1 & Data & 123 \\
Row 2 & More data & 456 \\
Row 3 & Even more & 789 \\
\hline
\end{tabular}
\end{table}

Reference tables like this: see Table \ref{tab:example}.

\section{Example of Equations}

Mathematical equations can be included inline like $E = mc^2$ or as display equations:

\begin{equation}
    \int_{a}^{b} f(x) \, dx = F(b) - F(a)
    \label{eq:fundamental}
\end{equation}

Reference equations using: see Equation \ref{eq:fundamental}.
